\documentclass[11pt]{article}
\usepackage{amsmath,amssymb}
\title{The connected first nonlinear response, including its original Haar partner}\author{}\date{}
\begin{document}\maketitle
\paragraph{Input and current scope.}
The mixed boundary magnetic loop bound calculation gives the full mixed-boundary loop proof.
Thus the complete signed leading constant-flux estimate is logarithmic,
not merely a conjecture from finite-size data. In the first original nonlinear flux coefficient calculation, the differentiated
fixed-$R$ Laplace expansion is justified: compactness, the unique tree-fixed
minimum, its positive Hessian and the fixed-dimensional action gap justify
its two parameter derivatives, but all constants may depend on $R$.
Its cubic-current test is not the complete nonlinear coefficient. Here
that coefficient is kept connected and evaluated exactly on the COMPLETE
original $R=2$ cube.

\paragraph{1. The connected formula before any size estimate.}
Use the original right-minimum tensors of the first original nonlinear flux coefficient calculation and its Gaussian
$\gamma_\epsilon$ with covariance $H_\epsilon^{-1}$. Write
\[
 T=T_{3,\epsilon},\quad W=T_{4,\epsilon}-J_2,\quad
 Y=W-\tfrac12 T^2,\quad
 q=\tfrac12 z^*H' z,\quad r=\tfrac12 z^*H''z,
 \qquad J_2=-\tfrac13\sum_j|z_j|^2.
\]
Then $\mathcal C_R=E_\gamma Y$. Differentiating the NORMALIZED Gaussian
expectation twice gives the exact finite-dimensional identity
\[
 \boxed{\mathcal C_R''=
 E Y''-2\operatorname{Cov}(Y',q)
       +\operatorname{cum}_3(Y,q,q)-\operatorname{Cov}(Y,r).}       \tag{1}
\]
Here
$\operatorname{cum}_3(Y,q,q)=E[(Y-EY)(q-Eq)^2]$.
Indeed $\partial E A=E A'-\operatorname{Cov}(A,q)$ and
$\partial\operatorname{Cov}(A,q)=\operatorname{Cov}(A',q)
+\operatorname{Cov}(A,r)-\operatorname{cum}_3(A,q,q)$.
All derivatives are in the actual right minimum frame; its velocity and
acceleration enter $H,T,W$. The Haar polynomial has no explicit parameter
derivative, but its expectation has covariance derivatives and cannot be
removed. Equation(1) removes disconnected vacuum products only through
actual cumulants; it does not prove their spatial summability.

On the complete $R=2$ cube every plaquette has at most two free links.
At zero background $T=0$ identically, so (1) becomes
\[
 \mathcal C_2''= E W''-E(T')^2
 -2\operatorname{Cov}(W',q)+\operatorname{cum}_3(W,q,q)
 -\operatorname{Cov}(W,r).                                        \tag{2}
\]
The cubic-source variance must be combined with all four other terms.
For larger cubes $T$ generally does not vanish, and (1), not (2), is required.

\paragraph{2. Exact complete-cube arithmetic.}
Take the physical constant-flux path of the complete cube abelian signed incidence calculation,
$a_{(x,\mu)}=\mathbf1_{\mu=1}x_0/4$, on $R=2$.
There are 216 original plaquettes, eight free links before gauge fixing,
and seven independent chords (21 real fluctuation coordinates).
The actual Cartan minimum is exactly this linear field for small physical
angle $\epsilon$. The tree edge is zero. All fixed-only plaquettes remain
in the action and the exact classical minimum.

The finite calculation expands each ORIGINAL ordered factor
\[
 U_e^*(\epsilon)\exp(i z_e\cdot\sigma)\quad(s_e=+1),\qquad
 \exp(-i z_e\cdot\sigma)U_e^*(\epsilon)^{-1}\quad(s_e=-1)
\]
to degree four in $z$ and degree two in $\epsilon$, using rational quaternion
products. It sums all plaquettes before taking Gaussian expectations.
The Hessian coefficients $H_0,H_1,H_2$ in
$H=H_0+\epsilon H_1+\epsilon^2H_2+\cdots$ give covariance coefficients
\[
 G_0=H_0^{-1},\quad G_1=-G_0H_1G_0,\quad
 G_2=-G_0(H_1G_1+H_2G_0).                                        \tag{3}
\]
Every inverse recurrence and covariance symmetry is checked exactly.
Wick pairing is used only to evaluate the finite Gaussian coefficient
already justified by the fixed-$R$ original Laplace expansion.
The twelve nonzero monomials of $T'_{3,0}$ are retained, not just a single
sample fluctuation.

The complete rational results are
\[
 \left(E_{\gamma_\epsilon}T_{4,\epsilon}\right)''_0
       =-{203\over5184},\qquad
 \left(-E_{\gamma_\epsilon}J_2\right)''_0={289\over6912},\qquad
 E_{\gamma_0}(T'_{3,0})^2={29\over3456}.
\]
Therefore
\[
 \boxed{\mathcal C_2''(0)=-{119\over20736},\qquad
        {\mathcal C_2''(0)\over E_2}=-{119\over46656},\quad
        E_2={9\over4}.}                                          \tag{4}
\]
The leading signed determinant computed by the same exact product jet is
$-71/576$, independently matching the complete cube abelian signed incidence calculation. A separate evaluation uses the
actual trigonometric minimum factors at nonzero $\epsilon$, expands only
in the fluctuations, and evaluates the entire Gaussian coefficient.
It gives $\mathcal C_2(0)=-3/8$ and centered second differences
\[
\begin{array}{c|r|r}
\epsilon&\text{second difference of }\mathcal C_2&\text{error from (4)}\\\hline
.04&-.00573868736045&1.244\,10^{-7}\\
.02&-.00573878063959&3.109\,10^{-8}\\
.01&-.00573880395538&7.773\,10^{-9}
\end{array}
\]
These checks verify the finite formula; they are not a proof of uniform
perturbation theory.

\paragraph{3. Consequence for the ORIGINAL fixed-block response.}
With $\epsilon=\zeta/\sqrt B$, the fixed-$R$ differentiated expansion gives
for this actual complete cube
\[
 F_{2,\zeta\zeta}(0)
   =-{71\over576B}-{119\over20736B^2}+O_{R=2}(B^{-3}).              \tag{5}
\]
Here $F=-\log Z-Bv$ is the explicitly classically-subtracted scalar.
The raw original $Z$ still has its full Haar/canonical normalization,
all fixed-only plaquette action, and the exact minimum $Bv$; those are not
reset or silently discarded. The complete tree Gaussian prefactor is
$(2\pi^2)^{-7}(2\pi/B)^{21/2}(\det H)^{-1/2}$ times $e^{-Bv}$.
There is no additional ghost factor in this already gauge-fixed tree
integral. Global modes remain outside the fixed relative cube.

\paragraph{What has and has not been closed.}
Equation(4) proves that the full first nonlinear response need not vanish
after its cubic, quartic and Haar partners are combined. It also shows why
the cubic variance alone predicts neither its size nor its sign: the
quartic and Haar responses largely cancel before that variance is subtracted.
The main remaining bound is on the connected combination(1), normalized
by physical boundary energy, together with a growing-$R$ differentiated
remainder or an actual-law resummation. The fixed-$R$ expansion(5) supplies
no such uniform control. The exact original-law residual (5) of the signed gauge quotient and unsummed current calculation remains the
nonperturbative formulation of that gap. No physical-scale iteration,
random-history admission or S94 conclusion follows from this coefficient.
\end{document}
